% 37
\begin{frame}{Можно ли распознать выпуклость через гессиан?}
\small
На лекции 2 мы получили формулу направленной кривизны:
\[
 \frac{d^2}{dt^2}f(x+tv)\Big|_{t=0}
 =v^\top\grad^2 f(x)v.
\]

Если функция выпукла, вдоль любой прямой она должна изгибаться \textbf{вверх}, а не вниз.

\centering
\includegraphics[width=0.72\linewidth]{hessian_geometries.pdf}
\end{frame}

% 38
\begin{frame}{Критерий второго порядка}
\small
\begin{block}{Теорема}
Пусть $f$ дважды непрерывно дифференцируема на выпуклом открытом множестве. Тогда
\[
\boxed{
 f\text{ выпукла}
 \iff
 \grad^2 f(x)\succeq0
 \quad\forall x.
}
\]
\end{block}

То есть выпуклость эквивалентна неотрицательной кривизне во всех направлениях и во всех точках.
\end{frame}

% 39
\begin{frame}{Почему выпуклость даёт $\grad^2 f\succeq0$}
\footnotesize
Зафиксируем точку $x$ и направление $v$. Рассмотрим одномерное сечение
\[
 \varphi(t)=f(x+tv).
\]
Если $f$ выпукла, то $\varphi$ тоже выпукла как функция одной переменной. Поэтому
\[
 \varphi''(t)\ge0.
\]
Но
\[
 \varphi''(t)
 =v^\top\grad^2f(x+tv)v.
\]
Значит, для любого $v$
\[
 v^\top\grad^2f(x)v\ge0,
\]
то есть $\grad^2f(x)\succeq0$.
\end{frame}

% 40
\begin{frame}{Обратное направление критерия второго порядка}
\footnotesize
Пусть теперь
\[
 \grad^2f(x)\succeq0
 \qquad\forall x.
\]
Для любых $x,y$ рассмотрим
\[
 \varphi(t)=f(x+t(y-x)).
\]
Тогда
\[
 \varphi''(t)
 =(y-x)^\top\grad^2f(x+t(y-x))(y-x)
 \ge0.
\]
Следовательно, $\varphi$ выпукла на $[0,1]$, а значит
\[
 \varphi(t)\le(1-t)\varphi(0)+t\varphi(1).
\]
Это ровно выпуклость $f$.
\end{frame}

% 41
\begin{frame}{Квадратичная функция как частный случай}
\small
Для
\[
 f(x)=\frac12x^\top Ax-b^\top x+c,
 \qquad A=A^\top,
\]
мы знаем
\[
 \grad^2f(x)=A.
\]
Поэтому
\[
\boxed{
 f\text{ выпукла}
 \iff
 A\succeq0.
}
\]

Это полностью согласуется со спектральной геометрией лекции 3.
\end{frame}

% 42
\begin{frame}{Практический алгоритм проверки выпуклости}
\small
Если функция дважды дифференцируема:
\[
\boxed{
 f
 \longrightarrow
 H(x)=\grad^2f(x)
 \longrightarrow
 \text{проверить }H(x)\succeq0\text{ для всех }x.
}
\]

\begin{block}{Но есть нюанс}
Если гессиан зависит от $x$, проверка положительной полуопределённости \textbf{во всех точках} может быть сама по себе нетривиальной задачей.
\end{block}
\end{frame}

% 43
\begin{frame}{Выпуклость не гарантирует единственность минимума}
\small
Рассмотрим
\[
 f(x,y)=x^2.
\]
Функция выпукла, но
\[
 f(0,y)=0
\]
для любого $y$.

То есть глобальных минимумов бесконечно много:
\[
 \argmin f=\{(0,y):y\in\R\}.
\]

Нужное усиление называется \textbf{строгой выпуклостью}.
\end{frame}

% 44
\begin{frame}{Строгая выпуклость}
\small
\begin{block}{Определение}
Функция $f$ называется \textbf{строго выпуклой}, если для любых $x\ne y$ и $t\in(0,1)$
\[
\boxed{
 f((1-t)x+ty)
 <
 (1-t)f(x)+tf(y).
}
\]
\end{block}

В отличие от обычной выпуклости, хорда теперь лежит \textbf{строго выше} графика между различными точками.
\end{frame}

% 45
\begin{frame}{Строгая выпуклость даёт единственность}
\small
\begin{block}{Теорема}
Строго выпуклая функция имеет не более одного глобального минимума.
\end{block}

Предположим, что $x\ne y$ --- два минимума и
\[
 f(x)=f(y)=f^\star.
\]
Тогда по строгой выпуклости
\[
 f\left(\frac{x+y}{2}\right)
 <
 \frac12f(x)+\frac12f(y)
 =f^\star,
\]
противоречие.
\end{frame}

% 46
\begin{frame}{Но строгая выпуклость не контролирует силу кривизны}
\centering
\includegraphics[width=0.83\linewidth]{strict_vs_strong.pdf}

\vspace{0.15em}
\small
Функция $x^4$ строго выпукла, но около нуля её вторая производная стремится к нулю. Значит, одной строгой выпуклости недостаточно для количественных оценок скорости.
\end{frame}

% 47
\begin{frame}{Сильная выпуклость}
\small
\begin{block}{Определение}
Дифференцируемая функция $f$ называется $\mu$-\textbf{сильно выпуклой}, если существует $\mu>0$ такое, что для любых $x,y$
\[
\boxed{
 f(y)\ge
 f(x)+\grad f(x)^\top(y-x)
 +\frac{\mu}{2}\norm{y-x}^2.
}
\]
\end{block}

По сравнению с обычной выпуклостью появляется дополнительный положительный квадратичный запас.
\end{frame}

% 48
\begin{frame}{Геометрия сильной выпуклости}
\centering
\includegraphics[width=0.81\linewidth]{strong_convexity_support.pdf}

\vspace{0.15em}
\small
Функция лежит выше не просто касательной, а \textbf{параболической опоры} с кривизной $\mu$.
\end{frame}

% 49
\begin{frame}{Критерий сильной выпуклости через гессиан}
\small
Для дважды дифференцируемой функции
\[
\boxed{
 f\text{ $\mu$-сильно выпукла}
 \iff
 \grad^2f(x)\succeq\mu I
 \quad\forall x.
}
\]

Удобный способ увидеть это: функция
\[
 g(x)=f(x)-\frac{\mu}{2}\norm{x}^2
\]
выпукла тогда и только тогда, когда
\[
 \grad^2g(x)=\grad^2f(x)-\mu I\succeq0.
\]
\end{frame}

% 50
\begin{frame}{Связь с гладкостью и числом обусловленности}
\small
Если одновременно
\[
 \mu I\preceq\grad^2f(x)\preceq LI
 \qquad\forall x,
\]
то кривизна функции во всех направлениях находится между $\mu$ и $L$.

Естественно определить
\[
\boxed{
 \kappa=\frac{L}{\mu}.
}
\]

\begin{alertblock}{Связь с лекциями 3–4}
Для квадратичной функции это ровно
\[
\kappa(A)=\frac{\lambda_{\max}(A)}{\lambda_{\min}(A)}.
\]
\end{alertblock}
\end{frame}
